\documentclass[UTF8,11pt]{ctexart}
\usepackage[a4paper,margin=24mm]{geometry}
\usepackage{amsmath,amssymb,amsthm,mathtools,mathrsfs}
\usepackage[all]{xy}
\usepackage{xurl,hyperref,enumitem}
\usepackage{tikz}
\usetikzlibrary{arrows.meta,cd,decorations.markings,calc,patterns}
\hypersetup{hidelinks,breaklinks=true}
\setlength{\emergencystretch}{3em}
\allowdisplaybreaks
\newtheorem*{theorem}{Theorem}
\newtheorem*{thm}{Theorem}
\newtheorem*{lemma}{Lemma}
\newtheorem*{proposition}{Proposition}
\newtheorem*{prop}{Proposition}
\newtheorem*{corollary}{Corollary}
\newtheorem*{cor}{Corollary}
\newtheorem*{claim}{Claim}
\theoremstyle{definition}
\newtheorem*{definition}{Definition}
\newtheorem*{defn}{Definition}
\newtheorem*{remark}{Remark}
\newtheorem*{example}{Example}
\newcommand{\R}{\mathbb R}
\newcommand{\C}{\mathbb C}
\newcommand{\Z}{\mathbb Z}
\newcommand{\Q}{\mathbb Q}
\newcommand{\N}{\mathbb N}
\newcommand{\F}{\mathbb F}
\newcommand{\D}{\mathbb D}
\newcommand{\bB}{\mathbb B}
\newcommand{\bC}{\mathbb C}
\newcommand{\bR}{\mathbb R}
\newcommand{\bZ}{\mathbb Z}
\newcommand{\bN}{\mathbb N}
\newcommand{\bQ}{\mathbb Q}
\newcommand{\bD}{\mathbb D}
\newcommand{\bF}{\mathbb F}
\newcommand{\CP}{\mathbb{CP}}
\newcommand{\RP}{\mathbb{RP}}
\newcommand{\sO}{\mathscr O}
\newcommand{\sI}{\mathscr I}
\newcommand{\sF}{\mathscr F}
\newcommand{\sG}{\mathscr G}
\newcommand{\sH}{\mathscr H}
\newcommand{\sR}{\mathscr R}
\newcommand{\sM}{\mathscr M}
\newcommand{\sV}{\mathscr V}
\newcommand{\sU}{\mathscr U}
\DeclareMathOperator{\Hom}{Hom}
\DeclareMathOperator{\Ext}{Ext}
\DeclareMathOperator{\Tor}{Tor}
\DeclareMathOperator{\Spec}{Spec}
\DeclareMathOperator{\Proj}{Proj}
\DeclareMathOperator{\supp}{supp}
\DeclareMathOperator{\im}{im}
\DeclareMathOperator{\id}{id}
\DeclareMathOperator{\Tr}{Tr}
\DeclareMathOperator{\tr}{tr}
\DeclareMathOperator{\rank}{rank}
\DeclareMathOperator{\ord}{ord}
\DeclareMathOperator{\dist}{dist}
\DeclareMathOperator{\End}{End}
\DeclareMathOperator{\Aut}{Aut}
\DeclareMathOperator{\Gal}{Gal}
\DeclareMathOperator{\vol}{vol}
\DeclareMathOperator{\Cl}{Cl}
\DeclareMathOperator{\coker}{coker}
\DeclareMathOperator{\codim}{codim}
\DeclareMathOperator{\divergence}{div}
\DeclareMathOperator{\Ric}{Ric}
\DeclareMathOperator{\grad}{grad}
\DeclareMathOperator{\Hess}{Hess}
\newcommand{\abs}[1]{\left\lvert#1\right\rvert}
\newcommand{\sabs}[1]{\lvert#1\rvert}
\newcommand{\babs}[1]{\bigl\lvert#1\bigr\rvert}
\newcommand{\Babs}[1]{\Bigl\lvert#1\Bigr\rvert}
\newcommand{\bbabs}[1]{\biggl\lvert#1\biggr\rvert}
\newcommand{\BBabs}[1]{\Biggl\lvert#1\Biggr\rvert}
\newcommand{\norm}[1]{\left\lVert#1\right\rVert}
\newcommand{\snorm}[1]{\lVert#1\rVert}
\newcommand{\bnorm}[1]{\bigl\lVert#1\bigr\rVert}
\newcommand{\Bnorm}[1]{\Bigl\lVert#1\Bigr\rVert}
\newcommand{\bbnorm}[1]{\biggl\lVert#1\biggr\rVert}
\newcommand{\BBnorm}[1]{\Biggl\lVert#1\Biggr\rVert}
\newcommand{\widebar}[1]{\overline{#1}}
\newcommand{\nicefrac}[2]{\frac{#1}{#2}}
\newcommand{\linnprod}[2]{\langle#1,#2\rangle}
\newcommand{\myindex}[1]{#1}
\newcommand{\glsadd}[1]{}
\newcommand{\avoidbreak}{}
\newcommand{\sourceref}[1]{\textnormal{[原文：\nolinkurl{#1}]}}
\newcommand{\thmref}[1]{\sourceref{#1}}
\newcommand{\lemmaref}[1]{\sourceref{#1}}
\newcommand{\propref}[1]{\sourceref{#1}}
\newcommand{\corref}[1]{\sourceref{#1}}
\newcommand{\defnref}[1]{\sourceref{#1}}
\newcommand{\exampleref}[1]{\sourceref{#1}}
\newcommand{\exerciseref}[1]{\sourceref{#1}}
\newcommand{\figureref}[1]{\sourceref{#1}}
\newcommand{\sectionref}[1]{\sourceref{#1}}
\newcommand{\appendixref}[1]{\sourceref{#1}}
\newcommand{\stacksref}[2]{\href{https://stacks.math.columbia.edu/tag/#1}{#2 [#1]}}


\newcommand{\E}{\mathbb E}
\newcommand{\Prob}{\mathbb P}
\newcommand{\Reff}{\mathcal R_{\mathrm{eff}}}
\newcommand{\eps}{\varepsilon}
\DeclareMathOperator{\diam}{diam}
\DeclareMathOperator{\Var}{Var}
\DeclareMathOperator{\Crit}{Crit}
\newcommand{\dd}{\,\mathrm d}


\begin{document}
\section*{082\quad 有界域Sobolev嵌入的Rellich紧致性}
\subsection*{来源与整理范围}
Benoit Dionne, \emph{Partial Differential Equations}, \texttt{sobolev.tex}，标签\texttt{sob\_rell\_th}，blob 9b4e4e05e38ca57ecd290f1fa38cc1f919c17ce4。
\par\url{https://github.com/BenoitDionne/PDE/blob/PDE/sobolev.tex}
\par\textbf{恢复方式（整理者说明）：}依据人类原文重新整理以补齐缺失题号；并非旧题证文件的逐字节恢复；2026-09-28。
\subsection*{作者命题与人类证明}

\paragraph{记号与假设（据作者章内约定整理）。}
$H^{k,2}=H^k$。$(k,2)$-extension property指存在有界线性延拓
$E:H^k(\Omega)\to H^k(\mathbb R^n)$，其限制回$\Omega$为恒等。
Fourier范数为$\|f\|_{k,\rho}^2=\int(1+|\xi|^2)^k|\widehat f(\xi)|^2\,d\xi$。
\begin{theorem}[Rellich's Theorem; sob\_rell\_th]
Let $0\leq k_1<k_2$ be integers. If $\Omega\subset\mathbb R^n$ is
bounded and has the $(k_2,2)$-extension property, the inclusion
$H^{k_2}(\Omega)\to H^{k_1}(\Omega)$ is compact.
\end{theorem}
\begin{proof}
Let $(f_i)$ be bounded in $H^{k_2}(\Omega)$.
Choose $\psi\in C_c^\infty(\mathbb R^n)$ equal to $1$ on $\Omega$ and
$0\leq\psi\leq1$. The operator $F=\psi E$ is bounded, satisfies
$F(f)|_\Omega=f$, and has $\supp F(f)\subset K:=\supp\psi$.
Put $g_i=F(f_i)$. Then
\[
\supp g_i\subset K,\qquad g_i|_\Omega=f_i,\qquad
\|g_i\|_{k_2,\rho}\leq C\quad\hbox{for all }i.
\]
\textbf{(i)} Choose $\phi\in C_c^\infty(\mathbb R^n)$ equal to $1$ on $K$.
Since $g_i=\phi g_i$, their Fourier transforms satisfy
$\widehat g_i=\widehat g_i*\widehat\phi$, with the normalization used for
this convolution identity. The convolution is smooth since
$\widehat\phi$ is Schwartz and $\widehat g_i\in L^2$.
The elementary weight inequality and Cauchy--Schwarz give
\begin{align*}
|\widehat g_i(y)|
&\leq\int|\widehat\phi(y-x)|\,|\widehat g_i(x)|\,dx\\
&\leq\frac{2^{k_2/2}}{(1+|y|^2)^{k_2/2}}
\int(1+|y-x|^2)^{k_2/2}|\widehat\phi(y-x)|
     (1+|x|^2)^{k_2/2}|\widehat g_i(x)|\,dx\\
&\leq2^{k_2/2}\|\phi\|_{k_2,\rho}\|g_i\|_{k_2,\rho}
\leq2^{k_2/2}C\|\phi\|_{k_2,\rho}.
\end{align*}
Differentiating the convolution, the same calculation, with
$\partial_j\widehat\phi$ in place of $\widehat\phi$, bounds
$|\partial_j\widehat g_i(y)|$ uniformly in $i,y$.
All its weighted $L^2$ norms are finite because $\widehat\phi$ is Schwartz.
Consequently the $\widehat g_i$ are uniformly bounded and equicontinuous,
the latter by the mean value theorem.

Arzel\`a--Ascoli gives a subsequence uniformly convergent on
$\overline{B_1(0)}$. Extract a further subsequence uniformly convergent
on $\overline{B_2(0)}$, and continue. The diagonal subsequence, denoted
$\widehat g_i^{[i]}$, converges uniformly on every $\overline{B_N(0)}$.

\textbf{(ii)} It suffices to prove that $(g_i^{[i]})$ converges in
$H^{k_1}(\mathbb R^n)$, because its restriction is $(f_i^{[i]})$.
Split the norm at radius $N$:
\begin{align*}
\|g_i^{[i]}-g_j^{[j]}\|_{k_1,\rho}^2
={}&\int_{|y|\leq N}(1+|y|^2)^{k_1}
 |\widehat g_i^{[i]}(y)-\widehat g_j^{[j]}(y)|^2\,dy\\
&+\int_{|y|>N}(1+|y|^2)^{k_1-k_2}(1+|y|^2)^{k_2}
 |\widehat g_i^{[i]}(y)-\widehat g_j^{[j]}(y)|^2\,dy\\
\leq{}&\sup_{|y|\leq N}|\widehat g_i^{[i]}(y)-\widehat g_j^{[j]}(y)|^2
 \int_{|y|\leq N}(1+|y|^2)^{k_1}\,dy\\
&+(1+N^2)^{k_1-k_2}\|g_i^{[i]}-g_j^{[j]}\|_{k_2,\rho}^2\\
\leq{}&\sup_{|y|\leq N}|\widehat g_i^{[i]}(y)-\widehat g_j^{[j]}(y)|^2
 \int_{|y|\leq N}(1+|y|^2)^{k_1}\,dy
 +4C^2(1+N^2)^{k_1-k_2}.
\end{align*}
Given $\epsilon>0$, first choose $N$ so that the last term is less
than $\epsilon/2$; this is possible because $k_1-k_2<0$.
For this fixed $N$, local uniform convergence supplies $J$ such that
the first term is less than $\epsilon/2$ for $i,j>J$.
Thus the diagonal sequence is Cauchy in the Banach space
$H^{k_1}(\mathbb R^n)$ and converges there. Restriction proves compactness.
\end{proof}

\subsection*{整理说明与先修范围（非作者正文）}
取作者末句和证明实际要求的严格阶差$k_1<k_2$；原定理首句的$\leq$不包含在本题题面。允许引用延拓算子（显式假设）、Sobolev乘子有界性、Plancherel和Arzel\`a--Ascoli。本题核心为公共紧支撑延拓、频率等度连续、对角子列及高低频分离，全数保留。原文对$\partial_j\widehat\phi$的范数写成$\partial_j\phi$，这里直接用前者的有限加权范数；末尾范数平方和常数按$\|g_i\|\leq C$统一为$4C^2$。这些是显示估计的记号校正。难度在于把弱的频率局部紧性提升为全局较低阶Sobolev范数收敛。
\subsection*{可修改的简短标注（非作者正文）}
$c$：函数空间嵌入的紧致性

$a$：有界延拓；紧支撑乘子；Fourier卷积；Arzelà--Ascoli；对角子列；频率尾部控制

\texttt{cross\_theory\_status = yes}。证明把空间域的公共支撑转为频率域的平滑与等度连续，再用导数阶差使频率尾部一致变小，返回Sobolev范数中的紧性。

全部标注为非穷尽、可修改初稿。
\subsection*{来源许可与编辑归属}
Benoit Dionne，\emph{Partial Differential Equations}，University of Ottawa，README署年2024，CC BY-NC-SA 4.0：\url{https://github.com/BenoitDionne/PDE/blob/PDE/README.md}；\url{https://creativecommons.org/licenses/by-nc-sa/4.0/}。本题编辑版本遵循同一许可。
ChatGPT 加入题号、中文来源说明、整理范围和初稿标注；编辑说明与人类证明分列。
\end{document}
